
\section{Conclusion and Discussion}

\revision{This study presents a privacy-preserving interpretable single-domain generalization framework for multi-label ocular disease recognition under restricted cross-institutional data access. DFE uses CAM guidance and dual-branch feature filtering to retain feature components that are relatively stable between original and style-augmented source images, and ARC models label-label and label-feature relevance while constraining the representation of each disease label to remain close across the two views. The single-source setting contributes to privacy: restricting updates to the accessible source domain avoids the target-domain access and cross-institutional data pooling that would otherwise be needed, and AdaGDP further protects the training information that remains in the source domain.}

\revision{The experimental results indicate that the gains come from the joint design. In both the ODIR Off-site and external EDID ablations, combining DFE and ARC outperforms either module alone. Compared with MixStyle, RSC, and FACT, the proposed method remains best on ODIR On-site and EDID in both the best-observed and mean candidate results, whereas FACT performs best on RFMiD. CAM visualizations and region-overlap analysis provide a way to inspect the image regions underlying the predictions. With AdaGDP, the selected model at a recorded $\epsilon$ of 9.76 ($\delta=10^{-5}$) maintains recognition performance close to the official non-private ODIR ResNet50 baseline.}

\revision{Several limitations remain. Although the asymmetric loss reduces the dominance of easy negative labels, class imbalance may still affect recognition stability for minority disease classes, and class-frequency-based weighting or resampling could be examined in future work. Cross-domain evaluation in this study relies on shared label spaces; future work can explore open-category and more complex out-of-distribution scenarios. For privacy preservation, future work can explore methods that retain performance under tighter privacy budgets, or combine the proposed mechanism with federated learning to support more complex cross-institutional deployment.}
